\chapter*{Аннотация}
\addcontentsline{toc}{chapter}{Аннотация}

     Измерена удельная теплоёмкость алюминия и титана с помощью калориметра. Удельная теплоёмкость титана оказалась равной $(550 \pm 40) \frac{\text{Дж}}{\text{К*кг}}$, алюминия - $(860.9 \pm 68) \frac{\text{Дж}}{\text{К*кг}}$  Предложенный метод измерения удельной теплоёмкости позволил экспериментально получить значения, которые в пределах погрешности совпадают с табличными.  Табличные значения удельных теплоёмкостей для алюминия и титана: $c_{\text{алюм}} = 902~\frac{\text{Дж}}{\text{кг*К}}$ $ c_{\text{титан}} = 530~\frac{\text{Дж}}{\text{кг*К}}$. Погрешность измерений для образца из титана составила 7\%, для образца из алюминия 8\%.

%и + этот метод позволяет определить теплоёмкость твёрдого тела с точностью, при комнатной температуре